\documentclass[11pt,a4paper]{article}
\usepackage[margin=1in]{geometry}
\usepackage[T1]{fontenc}
\usepackage{titlesec}
\usepackage{enumitem}
\usepackage{booktabs}
\usepackage{longtable}
\usepackage{xcolor}
\usepackage{fancyhdr}
\usepackage{hyperref}
\usepackage{parskip}

\hypersetup{
  colorlinks=true,
  linkcolor=blue!60!black,
  urlcolor=blue!60!black,
  citecolor=blue!60!black
}

\pagestyle{fancy}
\fancyhf{}
\fancyhead[L]{PetSocial --- Feature Specification}
\fancyfoot[C]{\thepage}
\renewcommand{\headrulewidth}{0.4pt}

\title{PetSocial\\\large Feature Specification}
\author{PetSocial Development Documentation}
\date{\today}

\begin{document}
\maketitle
\thispagestyle{fancy}
\tableofcontents

\section{Authentication}
\subsection{Sign up}
\begin{itemize}[leftmargin=*]
  \item Route: \texttt{/signup}. Fields: display name, username, email, password.
  \item Username is normalised to lowercase and checked for uniqueness against the
        \texttt{Owner} table before Supabase Auth account creation is attempted.
  \item If Supabase's email confirmation is enabled, the new account has no session and
        the user is redirected to \texttt{/signup/check-email}; if confirmation is
        disabled, the user is signed in immediately and redirected home.
  \item An \texttt{Owner} row is created in the same request, using the Supabase Auth
        user id as its primary key.
\end{itemize}

\subsection{Log in / Log out}
\begin{itemize}[leftmargin=*]
  \item Route: \texttt{/login}. Email and password only (no social login is currently
        wired up, though the UI groundwork for a Google button was scaffolded in an
        early design pass and never connected).
  \item Log out is a single Server Action that calls Supabase's \texttt{signOut()} and
        redirects to \texttt{/login}.
\end{itemize}

\section{Owner and Pet Profiles}
\subsection{Owner profile}
\begin{itemize}[leftmargin=*]
  \item Route: \texttt{/owners/[username]}.
  \item Displays display name (or ``Anonymous''), bio, follower/following counts, and a
        list of the owner's pets.
  \item Shows a \emph{Follow} button to any other signed-in visitor, and an
        \emph{Edit profile} button to the profile's own owner.
  \item Editing (\texttt{/owners/[username]/edit}) covers display name, bio, avatar
        upload, and the anonymity toggle.
\end{itemize}

\subsection{Pet profile}
\begin{itemize}[leftmargin=*]
  \item Route: \texttt{/pets/[id]}.
  \item Displays name, species, breed, bio, and the owner (or ``Managed by an anonymous
        owner'' if hidden).
  \item The pet's own owner additionally sees an \emph{Edit} button and, in a
        separate section on the same page, the pet's health log with an
        \emph{Add entry} action.
  \item Editing (\texttt{/pets/[id]/edit}) covers name, species, breed, bio, avatar
        upload, and the pet-level anonymity toggle.
\end{itemize}

\section{Posts and the Home Feed}
\begin{itemize}[leftmargin=*]
  \item Creating a post (\texttt{/posts/new}) requires a caption, a photo, or both. The
        post can optionally be attributed to one of the owner's own pets instead of the
        owner directly.
  \item The home feed (\texttt{/feed}) shows posts from the signed-in user and everyone
        they follow, newest first. A logged-out visitor sees a global feed of all posts
        instead.
  \item Each post supports a single like per owner (toggleable) and an unlimited number
        of comments.
  \item Post authorship display resolves through the anonymity rule described in the
        Database Design document: if the post is attributed to a pet, the pet's own
        \texttt{isAnonymous} flag \emph{or} the owner's flag hides the owner's name.
\end{itemize}

\section{Follows}
A follow is a directed edge from one owner to another. Following/unfollowing is a
single toggle action available from an owner's profile, the Explore page, and the
right-hand ``Who to follow'' panel on desktop. The home feed and (implicitly) the
notification stream are the only two features currently scoped by the follow graph.

\section{Communities}
\begin{itemize}[leftmargin=*]
  \item \texttt{/communities} lists every distinct species that has at least one pet,
        with a pet count for each.
  \item \texttt{/communities/[species]} lists that species' pets and its most recent
        posts, and offers a further filter by breed (via a query parameter, e.g.\
        \texttt{?breed=Corgi}).
  \item Communities are computed on the fly from existing \texttt{Pet.species} /
        \texttt{Pet.breed} values --- there is no separate ``community'' table to
        create or join.
\end{itemize}

\section{Lost \& Found}
\begin{itemize}[leftmargin=*]
  \item Route: \texttt{/lost-found} (list), \texttt{/lost-found/new} (report).
  \item A report has a status of either \texttt{lost} or \texttt{found}, a description,
        optional contact info, and an optional link to one of the reporter's own pets.
  \item The list page shows only unresolved reports. The reporting owner can mark their
        own report resolved, at which point it disappears from the active list (the row
        is kept, not deleted).
\end{itemize}

\section{Health Logs}
\begin{itemize}[leftmargin=*]
  \item Accessible only from a pet's own profile page, only to that pet's owner.
  \item An entry has a type (vaccination, vet visit, medication, or other), a title, an
        optional note, a date, and an optional reminder date.
  \item Health log data is never fetched at all when the viewer is not the pet's owner
        --- the visibility rule is enforced at the query level, not just hidden in the
        rendered page.
\end{itemize}

\section{Notifications}
\begin{itemize}[leftmargin=*]
  \item Route: \texttt{/notifications}.
  \item Generated for three events: someone liking your post, someone commenting on
        your post, and someone starting to follow you. An owner's own actions on their
        own content never generate a notification.
  \item The actor's name in a notification message respects that actor's own anonymity
        flag, showing ``Someone'' instead of their real name when appropriate.
  \item A ``Mark all read'' action is available whenever at least one notification is
        unread. The unread count also appears as a badge on the notification icon in
        the navigation.
\end{itemize}

\section{Direct Messages}
\begin{itemize}[leftmargin=*]
  \item Routes: \texttt{/messages} (conversation list), \texttt{/messages/[username]}
        (a single thread).
  \item The conversation list groups all of a user's messages by the other party,
        showing the most recent message and an unread count per conversation.
  \item Opening a thread marks every unread message \emph{from that specific sender} as
        read; messages in other conversations are unaffected.
  \item Direct messages are also the contact mechanism for both Adoption and
        Marketplace listings: a ``Message about \emph{item}'' button on a listing's
        detail page opens a thread with that listing's poster.
\end{itemize}

\section{Pet Adoption}
\begin{itemize}[leftmargin=*]
  \item Routes: \texttt{/adoption} (browse), \texttt{/adoption/new} (post a listing),
        \texttt{/adoption/[id]} (detail).
  \item A listing captures the pet's name, species (required), breed, age, gender,
        description (required), an optional photo, location, and contact info.
  \item The browse page can be filtered by species and always excludes listings whose
        status is \texttt{adopted}.
  \item The poster can cycle a listing's status between \texttt{available},
        \texttt{pending}, and \texttt{adopted} from the detail page; any other
        signed-in visitor sees a ``Message about \emph{name}'' button instead.
\end{itemize}

\section{Veterinary Booking}
\begin{itemize}[leftmargin=*]
  \item Routes: \texttt{/vet} (clinic directory), \texttt{/vet/[id]} (booking),
        \texttt{/vet/appointments} (the signed-in owner's own bookings).
  \item The clinic directory self-seeds with three demonstration clinics the first time
        it is visited on an empty \texttt{Vet} table; there is no clinic sign-up flow.
  \item The booking page shows the next seven days as date tabs. For the selected date
        it computes real availability from a fixed set of eight daily time slots
        (09:00--17:00, hourly except the 12:00 slot), excluding any slot that already
        has a confirmed appointment for that vet and excluding, for the current day
        only, any slot whose time has already passed.
  \item Booking requires selecting one of the signed-in owner's own pets and a reason
        (a fixed list: wellness check, vaccines, skin/coat, dental, other), plus
        optional notes.
  \item A booking is rejected server-side if the exact slot was already taken between
        page load and submission (a race-condition guard, not just a UI affordance).
  \item The owner can cancel their own upcoming appointment from
        \texttt{/vet/appointments}, which frees the slot for other bookings.
\end{itemize}

\section{Marketplace}
\begin{itemize}[leftmargin=*]
  \item Routes: \texttt{/marketplace} (browse), \texttt{/marketplace/new} (post a
        listing), \texttt{/marketplace/[id]} (detail).
  \item A listing has a title, category (food, toys, beds, health, accessories, other),
        price, description, an optional photo, and an optional location.
  \item The browse page can be filtered by category. Unlike an adoption listing, a sold
        marketplace listing remains visible in the browse list with a ``Sold'' badge
        rather than disappearing.
  \item The seller can toggle a listing between \texttt{available} and \texttt{sold}
        from the detail page; any other signed-in visitor sees a ``Message about
        \emph{title}'' button instead.
\end{itemize}

\section{Not Yet Implemented}
The following are either present in the data model but unused, or were explicitly
deferred during development and are not part of the current product:

\begin{itemize}[leftmargin=*]
  \item \textbf{Hashtags and text search.} The \texttt{Hashtag} and
        \texttt{HashtagsOnPosts} tables exist in the schema but no code path writes to
        or reads from them; there is no hashtag extraction from captions and no
        hashtag-based discovery page. The only existing search is a simple
        name/username filter on the Explore page.
  \item \textbf{Stories.} An ephemeral, 24-hour photo/video format shown at the top of
        the feed, as seen in comparable products. Not built.
  \item \textbf{Reels.} A short-form, full-screen vertical video feed. Not built.
  \item \textbf{Real-time chat delivery.} Direct messages are currently delivered by
        page revalidation (the recipient sees a new message on their next navigation or
        refresh, not via a live socket connection).
\end{itemize}

\end{document}
